\useoutertheme{infolines}
\usepackage{tikz}
\usepackage{amsmath}
\usepackage{amssymb}
\usepackage{mathrsfs}
\usepackage{isomath}
\usepackage{algorithmic}
\usepackage{fontawesome}
\usepackage{xcolor}
\usepackage{bbold}
\DeclareFontFamily{OT1}{mathc}{}
\DeclareFontShape{OT1}{mathc}{m}{it}{<-> mathc10}{}
\DeclareSymbolFont{mathc}{OT1}{mathc}{m}{it}

\DeclareMathSymbol{\mcalGamma}{\mathord}{mathc}{'203}
\DeclareMathSymbol{\mcalTheta}{\mathord}{mathc}{'210}
\DeclareMathSymbol{\mcalLambda}{\mathord}{mathc}{'213}
\DeclareMathSymbol{\mcalPi}{\mathord}{mathc}{'220}
\DeclareMathSymbol{\mcalTheta}{\mathord}{mathc}{'210}
\DeclareMathSymbol{\mcalalpha}{\mathord}{mathc}{'241}
\DeclareMathSymbol{\mcalbeta}{\mathord}{mathc}{'242}
\DeclareMathSymbol{\mcalgamma}{\mathord}{mathc}{'243}
\DeclareMathSymbol{\mcalvarepsilon}{\mathord}{mathc}{'245}
\DeclareMathSymbol{\mcaleta}{\mathord}{mathc}{'247}
\DeclareMathSymbol{\mcaltheta}{\mathord}{mathc}{'250}
\DeclareMathSymbol{\mcalkappa}{\mathord}{mathc}{'251}
\DeclareMathSymbol{\mcalmu}{\mathord}{mathc}{'253}
\DeclareMathSymbol{\mcalnu}{\mathord}{mathc}{'254}
\DeclareMathSymbol{\mcalomicron}{\mathord}{mathc}{'257}
\DeclareMathSymbol{\mcalpi}{\mathord}{mathc}{'260}
\DeclareMathSymbol{\mcalrho}{\mathord}{mathc}{'261}
\DeclareMathSymbol{\mcalupsilon}{\mathord}{mathc}{'264}
\DeclareMathSymbol{\mcalphi}{\mathord}{mathc}{'265}
\DeclareMathSymbol{\mcalchi}{\mathord}{mathc}{'266}
\DeclareMathSymbol{\mcalpsi}{\mathord}{mathc}{'267}
\DeclareMathSymbol{\mcalvartheta}{\mathord}{mathc}{'272}
\DeclareMathSymbol{\mcalvarrho}{\mathord}{mathc}{'275}
\DeclareMathSymbol{\mcalvarphi}{\mathord}{mathc}{'277}

\newcommand \E {\mathop{\mbox{\ensuremath{\mathbb{E}}}}\nolimits}
\newcommand \Var {\mathop{\mbox{\ensuremath{\mathbb{V}}}}\nolimits}
\newcommand \Bias {\mathop{\mbox{\ensuremath{\mathbb{B}}}}\nolimits}
\newcommand\ind[1]{\mathop{\mbox{\ensuremath{\mathbb{I}}}}\left\{#1\right\}}
\newcommand{\indep}{\mathrel{\text{\scalebox{1.07}{$\perp\mkern-10mu\perp$}}}}
\renewcommand \Pr {\mathop{\mbox{\ensuremath{\mathbb{P}}}}\nolimits}
\DeclareMathOperator*{\argmax}{arg\,max}
\DeclareMathOperator*{\argmin}{arg\,min}
\DeclareMathOperator*{\sgn}{sgn}
\newcommand \defn {\mathrel{\triangleq}}
\newcommand \loss {L}
\newcommand \step {\alpha}
\newcommand \cost {\ell}
\newcommand \util {U}
\newcommand \rewrad {r}
\newcommand \Reals {\mathbb{R}}
\newcommand \Naturals {\mathbb{N}}
\newcommand \pol {\pi}
\newcommand \Pols {\Pi}
\newcommand \Param {\mcalTheta}
\newcommand \param {\theta}
\newcommand \vparam {\vectorsym{\theta}}
\newcommand \mparam {\matrixsym{\Theta}}
\newcommand \CD {\mathcal{D}}
\newcommand \CX {\mathcal{X}}
\newcommand \CA {\mathcal{A}}
\newcommand \CY {\mathcal{Y}}
\newcommand \bW {\matrixsym{W}}
\newcommand \bw {\vectorsym{w}}
\newcommand \wi {\vectorsym{w}_i}
\newcommand \wij {w_{i,j}}
\newcommand \bA {\matrixsym{A}}
\newcommand \ai {\vectorsym{a}_i}
\newcommand \aij {a_{i,j}}
\newcommand \by {\vectorsym{y}}
\newcommand \hby {\vectorsym{\hat{y}}}
\newcommand \hy {\hat{y}}
\newcommand \bz {\vectorsym{z}}
\newcommand \bx {\vectorsym{x}}
\newcommand \bxt {\vectorsym{x}^t}
\newcommand \xti {x_i^t}
\newcommand \xtj {x_j^t}
\newcommand \xt {x^t}
\newcommand \yt {y^t}
\newcommand \bX {\matrixsym{X}}
\newcommand \bI {\matrixsym{I}}
\newcommand \dbel {\phi}
\newcommand \bel {\phi}
\newcommand \Bel {\Phi}
\newcommand \Ber {\textrm{Bernoulli}}
\newcommand \Unif {\textrm{Unif}}
\newcommand \Mult {\textrm{Mult}}
\newcommand \Beta {\textrm{Beta}}
\newcommand \Dir {\textrm{Dir}}
\newcommand \Wish {\textrm{Wish}}
\newcommand \Normal {\textrm{Normal}}
\newcommand \Simplex {\mathbb{\Delta}}
\newcommand \outcome {\omega}
\newcommand \Outcomes {\Omega}
\newcommand \cset[2] {\left\{#1 : #2 \right\}}
\usetikzlibrary{shapes.geometric}
\tikzstyle{utility}=[diamond,draw=black,draw=blue!50,fill=blue!10,inner sep=0mm, minimum size=8mm]
\tikzstyle{select}=[rectangle,draw=black,draw=blue!50,fill=blue!10,inner sep=0mm, minimum size=6mm]
\tikzstyle{hidden}=[dashed,draw=black,fill=red!10]
\tikzstyle{RV}=[circle,draw=black,draw=blue!50,fill=blue!10,inner sep=0mm, minimum size=6mm]
\tikzstyle{transition}=[rectangle,draw=black!50,fill=black!20,thick]
\tikzstyle{someset}=[circle,draw=black,minimum size=8mm]
\tikzstyle{point}=[circle,draw=black,fill=black]

